\section{Work-product and vertical-system inventory}\label{sec:C24:catalog}
Table~\ref{tab:C24:catalog} completes the catalog with research, business, vertical, and memory systems. It describes the role established by the inspected resource and the evidence needed before accepting a stronger claim. Product families can overlap; the categories organize experimental questions rather than define exclusive market segments. Names and source redirects are preserved without importing prior rankings or unverified ownership assertions.
\begingroup\setlength{\tabcolsep}{3pt}\small
\bookneedspace{.42\textheight}
\begin{longtable}{L{25mm}L{35mm}L{48mm}}
\caption{Research enterprise and memory products with role-specific checks.}\label{tab:C24:catalog}\\
\toprule Product & Inspected role & Diagnostic condition \\
\midrule\endfirsthead
\toprule Product & Inspected role & Diagnostic condition \\
\midrule\endhead
Mem0\cite{executionproduct024} & Fact-oriented memory layer & Test corrections, source provenance and deletion across derived records. \\
DeerFlow\cite{executionproduct029} & Inspectable research harness & Pin search tools and measure contradictory-evidence coverage at matched budget. \\
Graphiti / Zep\cite{executionproduct033} & Temporal graph context & Distinguish event time, knowledge time and invalidated facts during retrieval. \\
watsonx Orchestrate\cite{executionproduct034} & Enterprise orchestration & Verify connector principal propagation and business-process constraints at dispatch. \\
Hindsight\cite{executionproduct040} & Retention recall and reflection & Audit whether reflection adds unsupported inferences or preserves source links. \\
Agentforce\cite{executionproduct045} & CRM-grounded business agents & Change account state mid-run and test action preconditions and approval binding. \\
AI Agent Studio / Otto AI Agents\cite{executionproduct046} & Service-workflow agent tooling & Distinguish workflow completion from successful resolution and absence of adverse effects. \\
ChatGPT Work\cite{executionproduct047} & Research and connected work artifacts & Audit claim support and the exact saved artifact; hosted internals remain opaque. \\
Claude: work and Research\cite{executionproduct048} & Research and file-oriented work & Separate merged product naming from evidence about session and permission semantics. \\
AI by Zapier\cite{executionproduct053} & Cross-application adaptive automation & Inspect migration behavior and whether connector grants remain purpose scoped. \\
Asta\cite{executionproduct058} & Scientific search and analysis & Distinguish public research baselines from the hosted implementation and its tool access. \\
Cortex Agents\cite{executionproduct059} & Structured and unstructured data agents & Verify generated query authority and source-level access enforcement. \\
Gemini Deep Research / Deep Research Max\cite{executionproduct063} & Hosted multi-step research & Record preview or deployment tier and audit source coverage rather than report length. \\
Glean Agents\cite{executionproduct064} & Company-knowledge and action agents & Revoke document access after ingestion and test retrieval and generated summaries. \\
Copilot Studio\cite{executionproduct072} & Low-code business agents & Bind human approvals to exact tool arguments and current identity. \\
Manus\cite{executionproduct076} & Sandboxed general work & Test artifact persistence and unresolved external effects after an interrupted task. \\
Perplexity Computer\cite{executionproduct077} & Research-to-work interfaces & Measure transition from cited conclusions to correctly executed deliverables. \\
Rovo Agents\cite{executionproduct078} & Workplace knowledge and workflow agents & Check project boundaries and whether inherited permissions survive connector changes. \\
North / Agent Studio\cite{executionproduct079} & Enterprise agent construction & Treat private deployment and data isolation as configuration-specific claims requiring evidence. \\
Notion Custom Agents\cite{executionproduct080} & Scheduled workspace workflows & Inject duplicate triggers and inspect idempotency, current permissions and retention. \\
Joule Studio\cite{executionproduct081} & Business skills and SAP context & Test schema evolution and process-state constraints rather than conversational plausibility. \\
MemOS\cite{executionproduct083} & Composable memory interfaces & Identify actual update targets and distinguish memory writes from parameter learning. \\
OpenViking\cite{executionproduct084} & Context filesystem and retrieval & Inspect license artifact and additional restrictions separately; test deletion propagation. \\
Dust\cite{executionproduct085} & Company-context custom agents & Audit tool scope and evidence provenance across connected knowledge stores. \\
Kimi Agent\cite{executionproduct088} & Research analysis and work artifacts & Verify calculations and source support independently of presentation quality. \\
Harvey Agents / Agent Builder\cite{executionproduct090} & Legal workflow agents & Test jurisdiction, document versions, confidentiality and reviewer detection of seeded errors. \\
Sierra\cite{executionproduct092} & Customer-service agents & Measure policy compliance, resolution and escalation together with harmful-action rates. \\
WRITER Agent\cite{executionproduct093} & Enterprise playbooks and governance & Test whether reusable procedures remain valid after policy and model changes. \\
Grok Bot\cite{executionproduct094} & Persistent work-computer surface & Inspect identity, cloud-state lifetime and cancellation of background effects. \\
Copilot Researcher\cite{executionproduct098} & Web and enterprise research & Separate multi-model agreement from independent source verification. \\
Qianfan Agent Platform\cite{executionproduct099} & Regional agent-platform identity lead & Detailed mechanisms unresolved after source retrieval failure; no capability inference. \\
Genspark Super Agent\cite{executionproduct100} & Hosted research and deliverable creation & Check citation support, computation provenance and the final exported artifact. \\
\bottomrule\end{longtable}
\endgroup


The unresolved Qianfan source is included as an identity lead with unknown detailed behavior. This is more informative than silently omitting a regional product or guessing its controls from a neighboring offering. The same principle applies to feature parity across regions, deployment tiers, and private-cloud options: absence of inspected evidence constrains the claim. It does not establish absence of the feature.

Legal and customer-service systems require domain-specific outcomes. A legal workflow should be evaluated for source validity, jurisdiction, document consistency, review burden, and unauthorized disclosure. A support agent should be evaluated for policy compliance, resolution, escalation, and adverse actions, not only conversational satisfaction. A general benchmark can supply useful stress tests, but it cannot substitute for the domain's operational specification.

